% --- Archivo: 03_opt_unidimensional.tex ---

% =========================================================
\section{Métodos de optimización unidimensional}
\label{sec:v2_opt_unidimensional}
% =========================================================

La optimización unidimensional (conocida comúnmente como \textit{búsqueda de línea} o \textit{line search}) es un componente estructural crítico en la gran mayoría de los métodos multivariables. Tras calcular una dirección de descenso $d^k \in \Rn$ desde el punto actual $x^k$, la determinación del tamaño de paso $\alpha$ requiere resolver (exacta o aproximadamente) el siguiente subproblema:
\begin{equation}
    \min f(x) \quad \text{sujeto a} \quad x \in [a, b],
    \label{eq:v2_unidimensional_problem}
\end{equation}
donde $a < b$ y $f : [a, b] \to \R$ es una función continua.

Los algoritmos que presentaremos en esta sección son de \textbf{orden cero}, lo que significa que no exigen el cálculo de la derivada de $f$. Iniciaremos con el método de rejilla (un enfoque pasivo) y luego avanzaremos hacia los métodos secuenciales, que permiten aproximar la solución con un costo computacional mucho menor.

% ---------------------------------------------------------
\subsection{Búsqueda en rejilla (Grid Search)}
% ---------------------------------------------------------

El enfoque pasivo más intuitivo consiste en evaluar la función objetivo sobre una discretización uniforme del intervalo de búsqueda.

\begin{algorithm}[Búsqueda en rejilla]\label{alg:v2_grid_search}
    Fijamos un número entero $N \ge 1$ de puntos de evaluación permitidos.
    \begin{enuthm}
        \item Definimos el tamaño del incremento (resolución) de la rejilla como $\delta_N = (b-a)/N$.
        \item Calculamos los puntos de muestreo distribuidos simétricamente en el intervalo:
        \begin{equation}
            x^N_i = a + i \delta_N - \frac{\delta_N}{2}, \quad i = 1, \dots, N.
        \end{equation}
        \item Determinamos la aproximación devolviendo el valor mínimo obtenido tras evaluar los $N$ puntos:
        \begin{equation}
            v_N = \min_{i = 1, \dots, N} f(x^N_i).
        \end{equation}
    \end{enuthm}
\end{algorithm}

Si la función $f$ es Lipschitz continua en el intervalo con constante de Lipschitz $L > 0$ (es decir, $\absolute{f(y) - f(z)} \le L \absolute{y - z}$ para todo $y, z \in [a,b]$), el error analítico de la aproximación está rigurosamente acotado. 

Dado que cualquier minimizador global exacto $\bar{x} \in [a, b]$ dista a lo sumo la mitad de la resolución de la rejilla ($\delta_N/2$) del nodo de evaluación más cercano $x^N_j$, el error funcional del método satisface:
\begin{equation}
    v_N - \bar{v} \le f(x^N_j) - f(\bar{x}) \le L \absolute{x^N_j - \bar{x}} \le \frac{L(b - a)}{2N}.
    \label{eq:v2_grid_search_error}
\end{equation}

\begin{notabox}[Características de la Búsqueda en Rejilla]
    \
    \begin{enumerate}[label=\normalfont\arabic*., leftmargin=5mm]
        \item \textbf{Precisión óptima del error:} La cota de error obtenida en \eqref{eq:v2_grid_search_error} es ajustada (\textit{sharp}). Para todo número entero $N$ existe una función $f$, Lipschitz continua con constante $L$, tal que el error real de aproximación es exactamente igual a la cota teórica. Esto demuestra que, dentro de la clase de métodos \textbf{pasivos} de orden cero, el algoritmo de rejilla es óptimo en el sentido minimax.

        \item \textbf{Dificultad de estimación práctica:} La principal limitación práctica para definir a priori el número $N$ que garantiza una precisión funcional $\epsilon$ consiste en estimar adecuadamente el valor de la constante de Lipschitz $L$. En la gran mayoría de los casos reales, esta tarea no es trivial y puede resultar analíticamente tan compleja como resolver el propio problema de optimización original. Por lo tanto, el \cref{alg:v2_grid_search} suele tener utilidad computacional principalmente cuando $L$ es conocido por la física del problema o puede estimarse con un costo muy bajo.
    \end{enumerate}
\end{notabox}


% ---------------------------------------------------------
\subsection{Métodos de exclusión de intervalos}
% ---------------------------------------------------------

Los métodos secuenciales operan bajo un principio diferente: reducen la longitud del intervalo de incertidumbre de forma iterativa y adaptativa. Su análisis y garantía de convergencia requiere habitualmente que la función sea \textbf{unimodal}; es decir, que posea un único valle.

\begin{definition}[Función unimodal]\label{def:v2_unimodal}
    Se dice que una función continua $f : [a, b] \to \R$ es \deftech{unimodal} en $[a, b]$ si posee un único minimizador global $\bar{x} \in [a, b]$ y satisface dos condiciones de monotonía:
    \begin{itemize}
        \item Es estrictamente decreciente en el intervalo izquierdo $[a, \bar{x}]$.
        \item Es estrictamente creciente en el intervalo derecho $[\bar{x}, b]$.
    \end{itemize}
\end{definition}

\begin{remark}
    De la \cref{def:v2_unimodal} se deducen propiedades fundamentales:
    \begin{enuthm}
        \item Una función continua en $[a, b]$ es unimodal si, y solamente si, posee un \textbf{único minimizador local} en todo el dominio (el cual forzosamente coincide con el mínimo global).
        \item Toda función continua que sea \textbf{fuertemente convexa} en el intervalo $[a, b]$ es automáticamente unimodal.
        \item La unimodalidad es un concepto topológico \textbf{más débil que la convexidad}. Existen múltiples funciones unimodales que son estrictamente no convexas. Por ello, los métodos de exclusión de intervalos que presentaremos tienen un campo de aplicación mucho más amplio que las técnicas basadas únicamente en Análisis Convexo.
    \end{enuthm}
\end{remark}

El fundamento analítico subyacente de los métodos de exclusión (como Bisección o la Sección Áurea) radica en que la simple comparación ordinal de la función evaluada en dos puntos interiores permite acotar y garantizar la ubicación del mínimo global en un subintervalo estrictamente menor:

\begin{lemma}[Exclusión de intervalos]\label{lem:v2_interval_exclusion}
    Sea $f : [a, b] \to \R$ una función unimodal con minimizador global en $\bar{x} \in [a, b]$. Para cualesquiera dos puntos de prueba interiores $y, z \in [a, b]$ tales que $y < z$, se verifica inexorablemente la siguiente alternativa lógica:
    \begin{enuthm}
        \item Si $f(y) \le f(z)$, entonces el mínimo global se encuentra contenido en $\bar{x} \in [a, z]$.
        \item Si $f(y) \ge f(z)$, entonces el mínimo global se encuentra contenido en $\bar{x} \in [y, b]$.
    \end{enuthm}
\end{lemma}

\begin{proof}
    Demostraremos el caso (1) procediendo por reducción al absurdo. 
    
    Supongamos que se cumple $f(y) \le f(z)$, pero asumamos falsamente que el mínimo global $\bar{x}$ se encuentra estrictamente a la derecha del punto de corte $z$ (es decir, $\bar{x} > z$). En este escenario geométrico, los puntos mantienen el orden estricto $y < z < \bar{x}$.
    
    Invocando la definición de unimodalidad (\cref{def:v2_unimodal}), sabemos que la función $f$ debe ser \textbf{estrictamente decreciente} en todo el segmento anterior al mínimo $[a, \bar{x}]$. Como los puntos $y$ y $z$ pertenecen a este segmento de descenso, la relación espacial $y < z$ implica forzosamente que el valor funcional desciende: $f(y) > f(z)$. 
    
    Esto contradice frontalmente nuestra premisa inicial ($f(y) \le f(z)$). Por consiguiente, la asunción inicial era falsa y el mínimo global $\bar{x}$ no puede estar a la derecha de $z$. Debe pertenecer al intervalo cerrado $[a, z]$.
    
    La demostración del caso (2) es completamente análoga, apelando a la condición de crecimiento estricto en el segmento posterior al mínimo.
\end{proof}

\subsubsection{Método de Bisección}

El método de bisección divide sistemáticamente el dominio de búsqueda a la mitad, requiriendo evaluar la función en nuevos puntos centrales.

\begin{algorithm}[Método de bisección]\label{alg:v2_bisection}
    Fijamos los bordes $a_1 = a$, $b_1 = b$. Calculamos el centro inicial $c_1 = (a + b)/2$ y evaluamos funcionalmente $f(c_1)$. Para cada iteración $k \ge 1$:
    \begin{enuthm}
        \item Definimos el punto de prueba en el subintervalo izquierdo $y_k = (a_k + c_k)/2$ y evaluamos $f(y_k)$.
        \item Si se cumple $f(y_k) \le f(c_k)$, el \cref{lem:v2_interval_exclusion} garantiza que el mínimo está a la izquierda de $c_k$. El nuevo intervalo de incertidumbre es $[a_k, c_k]$. Actualizamos las variables:
        \[
            a_{k+1} = a_k, \quad b_{k+1} = c_k, \quad c_{k+1} = y_k.
        \]
        Saltamos directamente al paso de incremento (Paso 5).
        \item Si $f(y_k) > f(c_k)$, el mínimo debe estar a la derecha de $y_k$. Para decidir entre el bloque central y el derecho, definimos un punto de prueba en el subintervalo derecho $z_k = (c_k + b_k)/2$ y evaluamos $f(z_k)$.
        \item Si $f(c_k) \le f(z_k)$, actualizamos con el bloque central:
        \[
            a_{k+1} = y_k, \quad b_{k+1} = z_k, \quad c_{k+1} = c_k.
        \]
        Si $f(c_k) > f(z_k)$, actualizamos con el bloque derecho:
        \[
            a_{k+1} = c_k, \quad b_{k+1} = b_k, \quad c_{k+1} = z_k.
        \]
        \item Incrementamos el contador de iteraciones $k \leftarrow k+1$ y volvemos al Paso 1 hasta satisfacer una tolerancia $\epsilon$.
    \end{enuthm}
\end{algorithm}

La longitud del intervalo de incertidumbre se reduce a la mitad en cada paso ($b_{k+1} - a_{k+1} = (b_k - a_k)/2$). Sin embargo, la cantidad de evaluaciones de función requeridas es variable (a veces 1, a veces 2 por iteración). 

Si suponemos, por conveniencia algorítmica, que el presupuesto total de evaluaciones permitidas $N$ es un número impar, podemos realizar exactamente $k = (N-1)/2$ iteraciones completas. La cota de reducción del intervalo final sigue una progresión geométrica dictada por el factor $(1/2)^{1/2} \approx 0.707$:
\begin{equation}
    \text{Longitud Final} \le \frac{1}{2^{(N-1)/2}} (b - a) \approx (0.707)^{N-1}(b - a).
\end{equation}

\subsubsection{Método de la Sección Áurea}

Para maximizar la eficiencia computacional y garantizar que se requiera \textbf{exactamente una única evaluación de función} por cada iteración, se deben seleccionar los puntos de prueba interiores $y$ y $z$ manteniendo una proporción de invarianza geométrica constante $\tau$. 

Esta partición inteligente exige que la relación entre la longitud del intervalo original y el subintervalo mayor sea topológicamente idéntica a la relación entre el subintervalo mayor y el menor:
\begin{equation}
    \frac{b - a}{b - y} = \frac{b - y}{y - a}.
\end{equation}
Al resolver esta ecuación cuadrática de proporcionalidad, se obtiene que la constante de partición es el inverso de la proporción áurea (el límite de los cocientes de la sucesión de Fibonacci):
\begin{equation}
    \tau = \frac{\sqrt{5} - 1}{2} \approx 0.618.
\end{equation}

Para garantizar la validez del reciclaje de puntos internos sin necesidad de recalcularlos en cada paso, el \cref{alg:v2_golden_section} se apoya en el siguiente lema geométrico.

\begin{lemma}[Simetría e Invarianza de la Sección Áurea]\label{lem:v2_golden_properties}
    Sean $y$ y $z$ el menor y el mayor punto interior de la partición áurea del intervalo $[a, b]$, respectivamente. Entonces se verifica:
    \begin{enuthm}
        \item La distancia a los bordes es simétrica: $z - a = b - y = \tau(b - a)$.
        \item La propiedad de anidamiento fractal: el punto $y$ es exactamente el mayor punto interno de la partición de la sección áurea del subintervalo izquierdo $[a, z]$. Del mismo modo, $z$ es exactamente el menor punto interno de la partición áurea del subintervalo derecho $[y, b]$.
    \end{enuthm}
\end{lemma}

\begin{proof}
    La demostración exige una manipulación algebraica directa basada en las definiciones de $y$ y $z$, y se propone como el Ejercicio \ref{ex:v2_lema-seccion-aurea} al final del capítulo.
\end{proof}

Los puntos de prueba áureos para cualquier intervalo $[a, b]$ se calculan de forma exacta como:
\begin{align}
    y(a, b) &= a + \frac{3 - \sqrt{5}}{2}(b - a) \approx a + 0.382(b - a), \label{eq:v2_golden_y} \\
    z(a, b) &= a + \frac{\sqrt{5} - 1}{2}(b - a) \approx a + 0.618(b - a). \label{eq:v2_golden_z}
\end{align}

\begin{algorithm}[Método de la sección áurea]\label{alg:v2_golden_section}
    Fijamos los bordes $a_1 = a$, $b_1 = b$. Calculamos los dos puntos internos iniciales $y_1 = y(a_1, b_1)$ y $z_1 = z(a_1, b_1)$ mediante \eqref{eq:v2_golden_y}-\eqref{eq:v2_golden_z} y realizamos la primera evaluación dual $f(y_1)$ y $f(z_1)$. Tomamos el contador $k := 1$.
    \begin{enuthm}
        \item Para las iteraciones subsecuentes, \textbf{evaluamos únicamente} aquel punto interno (entre $y_k$ y $z_k$) cuyo valor de la función aún no haya sido calculado previamente.
        \item Invocamos el \cref{lem:v2_interval_exclusion}. Si $f(y_k) \le f(z_k)$, el mínimo está en $[a_k, z_k]$. Actualizamos los límites y reciclamos el punto interno evaluado (invocando el \cref{lem:v2_golden_properties}):
        \[
            a_{k+1} = a_k, \quad b_{k+1} = z_k, \quad z_{k+1} = y_k.
        \]
        Calculamos únicamente la coordenada del nuevo punto menor $y_{k+1} = y(a_{k+1}, b_{k+1})$.
        \item Si $f(y_k) > f(z_k)$, el mínimo está en $[y_k, b_k]$. Actualizamos y reciclamos en sentido inverso:
        \[
            a_{k+1} = y_k, \quad b_{k+1} = b_k, \quad y_{k+1} = z_k.
        \]
        Calculamos únicamente la coordenada del nuevo punto mayor $z_{k+1} = z(a_{k+1}, b_{k+1})$.
        \item Incrementamos $k \leftarrow k+1$ y volvemos al Paso 1.
    \end{enuthm}
\end{algorithm}

Al reciclar inteligentemente las evaluaciones funcionales, el factor de reducción del intervalo de incertidumbre por cada nueva función calculada es exactamente $\tau \approx 0.618$. El error final absoluto tras ejecutar un presupuesto de $N$ evaluaciones decae de manera geométrica:
\begin{equation}
    |x^N - \bar{x}| \le b_{k+1} - a_{k+1} = \left(\frac{\sqrt{5}-1}{2}\right)^{N-1}(b - a) \approx 0.618^{N-1}(b - a).
\end{equation}

% ---------------------------------------------------------
\subsection{Búsquedas en Dominios Semi-infinitos e Intervalos No Unimodales}
% ---------------------------------------------------------

\subsubsection{Fase de acotamiento en intervalos semi-infinitos}

En las implementaciones reales de métodos de gradiente, rara vez se conoce a priori un límite superior para la búsqueda de línea a lo largo de una dirección de descenso. Se requiere minimizar la función $f$ sobre una semirrecta unidimensional semi-infinita $[a, +\infty)$. Para poder instanciar el método de bisección o la sección áurea, la primera fase obliga a aislar topológicamente el mínimo en un intervalo acotado estricto $[a, a + T]$.

Esta \textbf{fase de acotamiento (bracketing)} funciona así: fijamos un tamaño de paso inicial exploratorio $\delta > 0$ y evaluamos la función en la secuencia de puntos equidistantes $z_k = a + k\delta$ para $k = 0, 1, 2, \dots$. 

Determinamos el menor índice entero $k \ge 0$ en el cual se viole la condición de descenso funcional, es decir:
\begin{equation}
    f(a + k\delta) \le f(a + (k+1)\delta).
\end{equation}
Bajo la hipótesis analítica de que $f$ mantiene su propiedad de unimodalidad a lo largo de toda la semirrecta, el mínimo global $\bar{x}$ pertenecerá invariablemente a la ventana acotada detectada:
\begin{equation}
    \bar{x} \in \begin{cases}
      [a, \, a + \delta] & \text{si } k = 0, \\
      [a + (k-1)\delta, \, a + (k+1)\delta] & \text{si } k \ge 1.
    \end{cases}
\end{equation}
Si el algoritmo fuera incapaz de encontrar dicho índice $k$ tras sucesivas evaluaciones (es decir, $f$ decrece asintóticamente al infinito), se concluiría matemáticamente que el problema unidimensional original es ilimitado y carece de minimizador finito.

\subsubsection{Robustez ante funciones no unimodales}

Un teorema pragmático de gran relevancia computacional dicta que si la función objetivo $f$ es continua en el intervalo cerrado $[a, b]$ pero \textbf{no es unimodal} (tiene múltiples valles locales y crestas), los algoritmos de Bisección (\cref{alg:v2_bisection}) y Sección Áurea (\cref{alg:v2_golden_section}) no colapsan ni arrojan errores de división por cero. 

En este escenario general caótico, los algoritmos operarán de forma ciega pero robusta, garantizando asintóticamente la convergencia hacia \textbf{al menos un minimizador local} (el cual depende estrictamente de los parámetros iniciales de partición) en el subespacio unidimensional.

% ---------------------------------------------------------
\subsection{Interpolación polinómica cuadrática y salvaguardas}
% ---------------------------------------------------------

Si la función objetivo $f$ es suficientemente suave (dos veces continuamente diferenciable), se puede abandonar la reducción geométrica pasiva y aproximar la topología del valle localmente mediante un polinomio cuadrático $P(x)$, tomando el vértice exacto de dicha parábola como la siguiente aproximación del mínimo. 

Este método requiere evaluar y mantener en memoria una terna de puntos $a_k < b_k < c_k$ que certifiquen el valle. A esta propiedad se le denomina \textbf{condición de sándwich}:
\begin{equation}
    a_k < b_k < c_k, \quad f(a_k) > f(b_k), \quad f(b_k) < f(c_k).
    \label{eq:v2_sandwich_cond}
\end{equation}

\begin{algorithm}[Interpolación cuadrática sucesiva]\label{alg:v2_quadratic_interpolation}
    Partiendo de una terna de puntos iniciales $a_1 < b_1 < c_1$ que garantice el cumplimiento de la condición de sándwich \eqref{eq:v2_sandwich_cond}, definimos el iterador $k := 1$:
    \begin{enuthm}
        \item Calculamos algebraicamente la abscisa del vértice $x_k$ de la parábola única que interpola los tres puntos:
        \begin{equation}
            x_k = \frac{1}{2} \frac{f(a_k)(c_k^2 - b_k^2) + f(b_k)(a_k^2 - c_k^2) + f(c_k)(b_k^2 - a_k^2)}{f(a_k)(c_k - b_k) + f(b_k)(a_k - c_k) + f(c_k)(b_k - a_k)}.
            \label{eq:v2_quadratic_vertex}
        \end{equation}
        \item Evaluamos computacionalmente la función objetivo en la nueva abscisa del vértice, es decir, obtenemos $f(x_k)$.
        \item Determinamos la topología del nuevo valle y actualizamos la terna analizando el árbol de decisiones lógicas:
        \begin{itemize}
            \item \textbf{Rama Derecha Inferior:} Si $x_k > b_k$ y $f(x_k) < f(b_k)$:
            \[
                a_{k+1} = b_k, \quad b_{k+1} = x_k, \quad c_{k+1} = c_k.
            \]
            \item \textbf{Rama Derecha Superior:} Si $x_k > b_k$ y $f(x_k) \ge f(b_k)$:
            \[
                a_{k+1} = a_k, \quad b_{k+1} = b_k, \quad c_{k+1} = x_k.
            \]
            \item \textbf{Rama Izquierda Inferior:} Si $x_k < b_k$ y $f(x_k) < f(b_k)$:
            \[
                a_{k+1} = a_k, \quad b_{k+1} = x_k, \quad c_{k+1} = b_k.
            \]
            \item \textbf{Rama Izquierda Superior:} Si $x_k < b_k$ y $f(x_k) \ge f(b_k)$:
            \[
                a_{k+1} = x_k, \quad b_{k+1} = b_k, \quad c_{k+1} = c_k.
            \]
        \end{itemize}
        \item Incrementamos $k \leftarrow k+1$ y volvemos al Paso 1.
    \end{enuthm}
\end{algorithm}

Bajo condiciones de regularidad local severas, el \cref{alg:v2_quadratic_interpolation} alcanza asintóticamente una tasa de convergencia \textbf{superlineal} de orden $\approx 1.32$ (mucho más rápida que la constante de la Sección Áurea). 

\begin{warningbox}[El colapso de la interpolación cuadrática]
    A diferencia de los métodos pasivos de exclusión de intervalos (que están ciegamente garantizados), la evaluación iterativa del denominador en la ecuación \eqref{eq:v2_quadratic_vertex} introduce una falla catastrófica potencial.
    
    Si los puntos de evaluación se alinean casi colinealmente, la curvatura de la parábola tiende a cero y el denominador experimenta un desbordamiento numérico. El vértice $x_k$ predicho puede ser arrojado violentamente fuera del intervalo de interés original. 
    
    Para evitar el colapso numérico, el conocido \textbf{Método de Brent} combina lo mejor de ambos mundos: intenta realizar un paso de interpolación cuadrática rápida, pero mantiene un estricto monitoreo. Si el paso cuadrático no resulta en una contracción suficiente o el vértice explota fuera de los límites, el algoritmo ejecuta una \textit{salvaguarda} recurriendo a un paso ciego y seguro del método de la Sección Áurea.
\end{warningbox}